%\input ../util/bookmacros
%\input ../util/docparams

\chapter{The Derivative As Approximation}

Here is the kind of question this chapter answers: suppose you are pouring a
concrete wall and your measurement of its height is off by an inch -- how far off
is the volume of concrete you computed? The tool that answers this turns out to be
the derivative, viewed in a new way. But before we get to the wall, let us warm up
by approximating integrals.

Sometimes, in fact many times, it may not be clear or even possible to evaluate an integral. It is
often the case that one can at least approximate the integral by bounding the value from above and below.
For instance, if $f$ is a continuous function%
\footnote{\kern 0.5pt \raise 0.5ex \hbox{\dag}}%
{Continuous functions are discussed in the section ``Computational Strategy for Limits'' in the appendix ``Infinite Sequences and Limits''.
They have the property that when drawing their graph, one never need take the pen off the page.} 
and non-negative
on an interval $[a,b]$, then $f$ attains a maximum, say $M$, and a minimum, say $m$ on this interval. Therefore,
the area under the graph of $f$ (the integral $\int_a^b f(x) \, dx$) 
is bounded above by the area of the rectangle with height $M$ and
width $b-a$ and bounded below by the area of the rectangle with height $m$ and width $b-a$.
That is, 
$$
m \, (b-a) \le \int_a^b f(x) \, dx \le M\, (b-a) 
$$
This result is true even if $f$ takes on negative values. This is true since the integral is 
a limit of sums of samples of $f$ along the interval, $[a,b]$,  times small differences -- which sum to $b-a$.
But each of these 
sample values of $f$ is bounded above and below by $M$ and $m$ respectively. 
Therefore, the sum is bounded above and below by $M \,(b-a)$ and $m \, (b-a)$.
We can show something just a little 
more interesting with almost no extra work.
\section{Mean Value Theorem}
Again, consider the area under the graph of 
$f$ from $a$ to $b$ ($a<b$)\footnote{\kern 0.5pt \raise 0.5ex \hbox{\ddag}}%
{The following argument works for non-negative $f$, but can be extended to all {\it continuous\/} $f$ (dropping the non-negativity assumption).}
: $\int_a^b f(s) \, ds$.
You can think of the graph as the outline of 
fine dust in a very narrow rectangular glass container. If we shift the dust around so that the dust
is level with height, say $h$, we must have (since the area of the dust is conserved)
$$
\eqalign {
\int_{a}^b f(s) \, ds & = h\,  (b - a) \cr
}
$$
If $m$ and $M$ are the minimum and maximum values that $f$ takes on 
in the interval $[a,b]$, then $m \le h \le M$. It turns out that continuous functions have 
the property that they attain all intermediate values. That is, if $f$ is a continuous 
function which hits the values $m$
and $M$ over the interval $[a,b]$, then for any $ m \le h \le M$ there exists a value 
$\mu \in [a,b]$ such that $f(\mu) = h$. Collecting these ideas we have:

\proclaim Theorem(\index{Intermediate Value Theorem}). If $f$ is a \index{continuous function} on the interval $[a, b]$ attaining values
$m$ and $M$ with $m \le M$, then for any \index{intermediate value} $h$ with $m \le h \le M$ 
there exists a value $\mu \in [a,b]$ such that $h = f(\mu)$.

\proclaim Theorem(\index{Mean Value Theorem}). If $f$ is a continuous function on the interval $[a,b]$,
then there exists a value, $\mu \in [a,b]$, such that
$$
\eqalign{
\int_{a}^b f(s) \, ds & = f(\mu) (b - a) \cr
}
$$

{\bf Note:\/} Many texts call this result the Mean Value Theorem {\it for
integrals\/}, reserving the name ``Mean Value Theorem'' for the tangent-line
statement, $f(b) - f(a) = f'(\mu)(b - a)$, derived from it below.

If $f'$ is also continuous then the theorem applies to $f'$:
$$
\int_a^b f'(s) \, ds = f'(\mu) (b - a)
$$
By the \index{Fundamental Theorem of Calculus} this can be written
$$
f(b) - f(a) = f'(\mu)(b - a)
$$
which is
$$
\eqalignno{
f(b) & = f(a) + f'(\mu)(b - a) & (1) \cr
}
$$
This can be written in a way that has an immediate \index{geometric interpretation}:
$$
\eqalign{
f'(\mu) & = {f(b) - f(a) \over (b - a)}  \cr
}
$$
This says that when $f'$ is \index{continuous} there is a point $\mu \in [a,b]$ where the \index{slope} of the 
\index{tangent line} is the same as the slope 
of the line passing through the points $(a, f(a))$ and $(b, f(b))$.

Here is another way to view this. Consider the interval $[x-h, x+h]$ for some $h>0$.  Substituting $x+h$ for $b$ and $x$ for $a$ in (1) gives
$f(x+h) = f(x) + f'(\mu_1)h$ -- for some $\mu_1 \in [x, x+h]$. Next substituting $x$ for $b$ and $x-h$ for $a$ in (1) gives
$f(x) = f(x-h) + f'(\mu_2) h$, or $f(x-h) = f(x) - f'(\mu_2)h$ -- for some $\mu_2 \in [x-h,x]$. That is, we may write $f(x+h)$ as
$f(x+h) = f(x) + f'(\mu)h$ regardless of the sign of $h$, for some $\mu \in [x-|h|, x+|h|]$.
If $h$ is small, then $f'(\mu)$ is close to $f'(x)$; therefore, an approximation to $f(x+h)$ is
$$
\eqalignno{
f(x+h) & \approx f(x) + f'(x) h & (1a) \cr
}
$$

\section{Derivative as Best Linear Approximation}
Equation (1a) amounts to a restatement of the definition of a derivative of $f$ at $x$.
This equation says that we may approximate $f(x+h)$ with a straight line; that is
a {\it \index{rectilinear}\/} function -- also called an {\it \index{affine}\/} function
(most books say simply {\it linear\/}).
This can be thought of as starting with $f(x)$ and adding a {\it \index{Linear Approximation}}
that is a linear function of the change in $x$, $h$.

In what sense is the tangent line the {\it best\/} line? Among all lines passing
through the point $(x, f(x))$, it is the one whose error goes to zero faster than
$h$ itself. Recall from the Mean Value Theorem that
$f(x+h) - \bigl(f(x) + f'(x)h\bigr) = \bigl(f'(\mu) - f'(x)\bigr)h$, and the factor
$f'(\mu) - f'(x)$ itself shrinks to zero as $h$ does. Any other line through
$(x, f(x))$, say with slope $s \ne f'(x)$, has error approximately
$\bigl(f'(x) - s\bigr)h$ -- proportional to $h$ itself.

This turns out to be a way to describe the derivative in a way that is extensible
to ``higher order'' functions; e.g., so-called {\it \index{multi-variable}} and
{\it \index{vector-valued}} functions.
By this we mean functions of the form:
$f: R^n \rightarrow R$ and $g: R^n \rightarrow R^m$.
It is the same game on a bigger board. For $f: R \rightarrow R$ the derivative is a
single number -- a slope. For $f: R^n \rightarrow R$ it is a vector, called the
{\it gradient\/}; and for $g: R^n \rightarrow R^m$ it is a matrix, called the
{\it Jacobian\/}. In every setting the answer to the question ``what is the
derivative?''\ is the same: it is the linear map that best approximates the
change in the function.

Such functions come under the general heading of ``Advanced Calculus'',
and are beyond the scope of this book.

\section{Approximation Application}
Suppose that you needed to build a concrete wall 100 ft long and 5 ft high. Builders tell you to set the depth of the wall
based on the height of the wall by the  following formula: depth (in inches) = height (in inches) / 10. So, for a given height $x$ in inches the volume of concrete
needed is $V(x) = 100\,{\rm ft} \times x\,{\rm in} \times (x/10)\,{\rm in}$. Using the units of inches, the formula for $V$ (in cubic inches) in
terms of height, $x$, in
inches is $V(x) = 12 \times 100 x^2/10 = 120 x^2$. To pour a 5 ft wall we need (x = 60 inches) $V(60) = 432{,}000$ cubic inches of
concrete, which is approximately 9 1/4 cubic yards.
But there is always error in measurement. What if, due to the ground conditions, we are off by 1 inch when measuring the
height? We need to order extra concrete to be sure we have enough. How much more should we order to take care of this 1 inch
error? We don't want to order too much extra because we want to control our costs. We could simply compute $V(60 + 1)$ easily
enough -- this gives an upper bound on the amount of concrete to use.
However, we would like to estimate the error  for {\it any\/} small error in height.% 
\footnote{\kern 0.5pt \raise 0.5ex \hbox{\ddag}}{In our case, $V(60+h) - V(60)$ is the extra amount of concrete to use. 
However, in many similar problems the approximation formula (1a) 
is often simpler and provides more insight into how sensitive the original function is to {\it all\/} small errors at a given input.}

Equation (1a) tells us that if there is a small error in the height then  the 
approximate change in the volume is given by $V(x+h) \approx V(x) + V'(x) h$. Let's say that we are concerned with an error of
1 in when the height is 5 ft (60 inches). The magnitude of the change is $|V'(x)h|$. It is easy to see that $V'(x) = 240
x$.  When $x = 60$ we have
 $V'(60) = 14400$; if, in addition, $h=1$,  $|V'(x)h|$ becomes $14400$ cubic inches. This is 0.309 cubic yards.
Therefore, we should order an extra 1/3 of a yard of concrete.

\exercise{Use the tangent line approximation (1a) with $f(x) = x^5$ to estimate
$(2.01)^5$ without a calculator. The true value is $32.80804\ldots$ -- how far
off is the estimate?}

\exercise{In the concrete wall problem, suppose the height measurement could be
off by as much as $2$ inches instead of $1$. Estimate the extra concrete needed.
Is your answer exactly twice the $1$ inch answer? Compare with the exact extra
amount, $V(62) - V(60)$.}

\section{Taylor's Theorem}
We now describe a result that is more difficult but which is quite impressive%
\footnote{\kern 0.5pt \raise 0.5ex \hbox{\dag}}{This section stands by itself and may be skipped if you wish.}.
Let us first 
describe a notation for differentiation.
We have used $f'$ to denote the derivative of $f$. We will use the notation $f''$ to denote the 
derivative of the function $f'$. Similarly, we use the notation $f'''$ to denote the derivative 
of $f''$. In general, we use the notation $f^{(n)}(x)$ to denote 
the $n^{\rm th}$ derivative of $f$ evaluated at $x$. 

Now, if $f'$ is a continuous function we know 
by the mean value theorem that
$$
\eqalign{
f(b) & = f(a) + f'(\mu)(b - a) \cr
}
$$
But we could invoke the theorem for any $x \in [a,b]$ because $[a,x]$ is an interval where the 
theorem applies. In this case the conclusion of the theorem is
$$
\eqalignno{
f(x) & = f(a) + f'(\mu_1)(x - a) & (2) \cr
}
$$
Note that we changed $\mu$ to $\mu_1$ because the interval $[a,x]$ is different than the interval $[a,b]$.
All that the \index{Mean Value Theorem} says is that there is {\it some\/} value in $[a,x]$; it may not 
be the same one that applied on the interval $[a,b]$. However, for the time being, let us assume 
that it is the same; that is, we assume that $\mu_1 = \mu$.

If $f''$ is also continuous on $[a,b]$, 
then we may apply this formula to $f'$ (it's a function which satisfies the requirements of the theorem) giving
$$
\eqalign{
f'(x) & = f'(a) + f''(\mu_2)(x - a) \cr
}
$$
for some value $\mu_2$ which depends on the interval $[a,x]$. Why is this different from $\mu_1$?
Because $f'$ is a different function from $f$ so we should not expect $\mu_2$ to be the same as $\mu_1$.
Let us assume, however, that 
$\mu_2 = \mu_1$, which we are assuming is $\mu$.

Integrating this function over the interval $[a,x]$ we have
$$
\eqalign{
\int_a^x f'(s) \, ds & = \int_a^x f'(a)\, ds  + \int_a^x f''(\mu)(s - a) \, ds \cr
}
$$
Since the \index{integration} is \index{linear}, we may place the constants outside of the \index{integral}%
\footnote{\kern 0.5pt \raise 0.5ex \hbox{\ddag}}%
{Anything that does not depend on the \index{integration} variable, $s$, is a constant
with respect to the integration. Therefore, it may be pulled outside of the integral. }.
$$
\eqalign{
\int_a^x f'(s) \, ds & = f'(a) \int_a^x 1 \, ds  + f''(\mu) \int_a^x (s - a) \, ds \cr
}
$$
Using the \index{FTC} on the left hand and first term of the right hand side gives
$$
\eqalign{
f(x) - f(a) & = f'(a) (x-a)  + f''(\mu) \int_a^x (s - a) \, ds \cr
}
$$
As an \index{anti-derivative} of $(s-a)$ is ${(s-a)^2 \over 2}$, we have
$$
\eqalign{
f(x) & = f(a) + f'(a) (x-a)  + f''(\mu) \int_a^x \biggl({(s - a)^2 \over 2}\biggr)' \, ds 
}
$$
Using the \index{FTC} on the last term, we may write
$$
\eqalignno{
f(x) & = f(a) + f'(a) (x - a) + f''(\mu) {(x-a)^2 \over 2} & (3)\cr
}
$$
Now, just as we plugged $f'$ into equation (2), we can do the same with equation (3).
If the function $f^{(3)}$ is continuous, then plugging $f'$ into (3) gives
$$
\eqalign{
f'(x) &= f'(a) + f''(a) (x - a) + f^{(3)}(\mu_3) {(x-a)^2 \over 2}  \cr
}
$$
Again, we need a new ``$\mu$'' because we are using a new function, $f'$.
As before, we integrate this equation over the interval $[a,x]$ giving
$$
\eqalign{
f(x) - f(a) &= f'(a) (x - a) + f''(a) {(x - a)^2 \over 2} + \int_a^x f^{(3)}(\mu_3) {(s-a)^2 \over 2}\, ds \cr
}
$$
As an \index{anti-derivative} of $(s-a)^2/2$ is ${(s-a)^3 \over 3 * 2}$, the above becomes
$$
\eqalign{
f(x) - f(a) &= f'(a) (x - a) + f''(a) {(x - a)^2 \over 2} + f^{(3)}(\mu_3) \int_a^x {(s-a)^2 \over 2}\, ds \cr
}
$$
Which is
$$
\eqalignno{
f(x) &= f(a) + f'(a) (x - a) + f''(a) {(x - a)^2 \over 2} + f^{(3)}(\mu_3) {(x-a)^3 \over 3 * 2} & (4) \cr
}
$$
We can keep plugging $f'$ into these successively more complicated equations (equations (2), (3), and (4)) 
provided $f$ 
has enough derivatives and they are continuous. The pattern that emerges
is called Taylor's Theorem:
\proclaim Theorem(Taylor's Theorem). If $f$ is a function such that $f^{(n)}$ exists 
and is a continuous function 
on the interval $[a,b]$, then for any $x \in [a,b]$\footnote{\kern 0.5pt \raise 0.5ex \hbox{\dag}}{%
The notation $n!$ means the product of all the integers from $1$ to $n$.}
$$
\eqalign{
f(x) &= f(a) + f'(a) (x-a) + f^{(2)}(a) {(x-a)^2 \over 2!} + f^{(3)}(a) {(x-a)^3 \over 3!} \cr
& + \ldots + f^{(n-1)}(a) {(x-a)^{n-1} \over (n-1)!} + f^{(n)}(\mu(x)) {(x-a)^n \over n!} \cr
}
$$
where $\mu$ is a function such that $\mu(x) \in [a,x]$.

Notice what this theorem says: from the values of $f$ and its derivatives at
the {\it single\/} point $a$, we can rebuild $f$ across the whole interval.
This is the same game we have been playing since the discrete chapters:
reconstruct the whole from its rates of change at a single point.

Our derivations for Taylor's Theorem in the cases $n=2$ and $n=3$ are not correct 
since we assumed at each stage that the various ``$\mu$'' were the 
same. A mathematically correct derivation can be found in the appendix on Taylor's Theorem. The 
derivation has the same structure as our derivation above but uses a generalization 
of the \index{Mean Value Theorem} to prove the result. One can also show without much work that 
the formula is true when $x$ is less than $a$:

\proclaim Corollary. If $f$ is a function such that $f^{(n)}$ exists 
and is a continuous function 
on the interval $[a,b]$, with $c\in [a,b]$, 
then for any $x \in [a,b]$
$$
\eqalign{
f(x) &= f(c) + f'(c) (x-c) + f^{(2)}(c) {(x-c)^2 \over 2!} + f^{(3)}(c) {(x-c)^3 \over 3!} \cr
& + \ldots + f^{(n-1)}(c) {(x-c)^{n-1} \over (n-1)!} + f^{(n)}(\mu(x)) {(x-c)^n \over n!} \cr
}
$$
where $\mu$ is a function such that $\mu(x)$ is a value between $x$ and $c$.

Taylor's Theorem says that we may approximate a function $f$ with $n$ continuous
derivatives, near the point $c$, with a
{\it polynomial\/} of degree $n-1$. The error in the approximation is the term
$f^{(n)}(\mu(x)) {(x-c)^n \over n!}$.

\exercise{Use Taylor's Theorem with $c = 0$ and $n = 3$ to show that
$\cos(x) \approx 1 - x^2/2$ with an error of size at most $|x|^3 / 6$. How good
is the approximation when $|x| \le 0.1$?}

\section{Power Series Representation}
In summation notation Taylor's formula may be written:
$$
f(x) = f(a) + \sum_{k=1}^{n-1} f^{(k)}(a) {(x-a)^k \over k!} + f^{(n)}(\mu(x)) {(x-a)^n \over n!}
$$
If a function has an infinite number of derivatives and the last term in the above formula
goes to zero in a sufficiently ``well behaved''  
manner (which is beyond the scope of this book to describe) as $n$ goes to $\infty$, then we may write
$$
f(x) = f(a) + \sum_{n=1}^\infty f^{(n)}(a) {(x-a)^n \over n!}
$$
One function that has an infinite number of derivatives is the exponential function, $f(x) = e^x$.
We know that $f'(x) = e^x$, and consequently $f^{(n)}(x) = e^x$. It is also a function that 
is sufficiently ``well behaved''. Letting $a=0$ we may write\footnote{{\kern 0.5pt \raise 0.5ex \hbox{\ddag}}}%
{The notation $0!$ is defined to be $1$. In this context we also adopt the convention that $x^0 = 1$ for {\it all\/} $x$ -- including $x = 0$ -- so that the $n=0$ term of the series makes sense at $x = 0$.}
$$
e^x = e^0 + \sum_{n=1}^\infty e^0 {x^n \over n!} = \sum_{n=0}^\infty {x^n \over n!}
$$
This says that $e^x$ can be thought of as an ``infinite polynomial''. Many other such functions can
be represented in this way, for instance the trig functions, $\sin$ and $\cos$.
Such representations can be used to find (among other things) an approximation to a function, its derivative, or its integral.

Finally, notice that the value $e$ itself can be approximated for a given $N$ by 
$$
e = e^1 = \sum_{n=0}^\infty {1^n \over n!}  = \sum_{n=0}^\infty {1 \over n!} \approx 1 + 1 + {1 \over 2!} + {1 \over 3!} + \ldots + {1 \over N!}
$$
% \bye

